\rok{Новембар 2023.}{C}

Први практични тест из Алгоритама и структура података

\zadatak{1}{Брисање елемената из реда (Queue Dequeue)}
Написати методу која имплементира стандардан начин за брисање елемената из реда.

\textbf{Додатне напомене за решавање задатка:}
\begin{itemize}
    \item Алгоритам написати у оквиру закоментарисане методе \texttt{dequeue?()}.
    \item Поменута метода се налази у оквиру класе \texttt{Queue} која се налази у придруженом \textit{template}-у.
    \item Обезбедити код у \texttt{Main} методи за тестирање алгоритма.
\end{itemize}



\zadatak{2}{Да ли је једноструко повезана листа бројева Фибоначијев секвенца?}
Креирати функцију која ће проверавати да ли бројеви у једноструко повезаној листи одговарају Фибоначијевој секвенци. Фибоначијева секвенца је секвенца целих бројева код којих се сваки број добија сумом његова два претходника у секвенци. Алгоритам је неопходно написати са линеарном временском комплексношћу $O(n)$.

\textbf{Додатни захтеви за решавање задатка:}
\begin{itemize}
    \item Алгоритам писати у оквиру класе \texttt{LinkedList}.
    \item Захтеван потпис методе:
    \begin{Verbatim}[fontsize=\footnotesize, numbers=left, xleftmargin=10mm]
public bool IsLinkedListFibonacci()
    \end{Verbatim}
\end{itemize}

\textbf{Примери:}
\begin{itemize}
    \item \textbf{Улаз:} $0 \rightarrow 1 \rightarrow 1 \rightarrow 2 \rightarrow 3 \rightarrow 5 \rightarrow 8 \rightarrow 13 \rightarrow 21$ \\
          \textbf{Излаз:} true
    \item \textbf{Улаз:} $1 \rightarrow 1 \rightarrow 2 \rightarrow 3$ \\
          \textbf{Излаз:} false
    \item \textbf{Улаз:} null \\
          \textbf{Излаз:} false
\end{itemize}



\zadatak{3}{Проналажење поднизова чија сума одговара прослеђеној вредности}
Креирати алгоритам који из низа целих бројева проналази све поднизове дужине 3 чија сума је једнака прослеђеној вредности. Алгоритам је неопходно написати са линеарно-логаритамском временском комплексношћу $O(n \cdot \log n)$.

\textbf{Додатни захтеви за решавање задатка:}
\begin{itemize}
    \item Вредности у оквиру поднизова не смеју да се понављају – не прихватају се различите комбинације бројева.
    \item У реализацији задатка може се користити алгоритам MergeSort који је дат у шаблону.
    \item Захтеван потпис методе:
    \begin{Verbatim}[fontsize=\footnotesize, numbers=left, xleftmargin=10mm]
public static void FindThreeNumbers(long[] array, int size, int sum)
    \end{Verbatim}
\end{itemize}

\textbf{Примери:}
\begin{itemize}
    \item \textbf{Улаз 1:} [2,4,3,6,0,8,5], 7, 10 \\
          \textbf{Излаз 1:} \\
          Triplet is: 0 2 8 \\
          Triplet is: 0 4 6 \\
          Triplet is: 2 3 5
    \item \textbf{Улаз 2:} [1,2,7,6,0,3], 6, 9 \\
          \textbf{Излаз 2:} \\
          Triplet is: 0 2 7 \\
          Triplet is: 0 3 6 \\
          Triplet is: 1 2 6
\end{itemize}

\sablon{AISP2023_Test1_GrupaC}{2023/Novembar/GrupaC/AISP2023_Test1_GrupaC.linq}
